\scanpage{576}
% source: upper portion, continuation of proof, lines 1--7 and (14).

inverse image in $\mathcal{G}^{+\wedge}$. For condition
b), we are reduced to the case where $u_{0}\in\mathcal{M}_{0,3}^{\sim}[0]$
(conjugating if necessary by an element of
$\mathcal{G}_{0,3}^{+\wedge}\simeq\mathcal{G}^{+\wedge}/\!\pm1$),
in this case we have by construction of $\mathcal{N}_{1,1}^{\sim}$ from
$\mathcal{M}_{0,3}^{\sim}$ (cf. \uline{\(\overline{\mathrm{III}}\)}, \unsure{pages|pages} 665...)
\[
(14)\qquad u_{0}(\varepsilon_{0})=\varepsilon_{0}^{\mu}
\]
where $\mu\in\widehat{\mathbb{Z}}^{*}$ is the multiplier, OK.

% source: middle portion, converse argument, lines 8--17.
Conversely, let $u$ be an automorphism of $\mathcal{G}^{+\wedge}$ which
satisfies a),b), consider the automorphism $\widetilde{u}$ of
$\mathcal{G}^{+\wedge}/\!\pm1\simeq\mathcal{G}_{0,3}^{+\wedge}$
which it induces. We thus see
that $\widetilde{u}(\Pi_{0,3}^{\wedge})=\Pi_{0,3}^{\wedge}$, and that
$\widetilde{u}(\varepsilon_{0})$ is $\mathcal{G}_{0,3}^{\wedge}$-conjugate
to an $\varepsilon_{0}^{\mu}$ ($\mu\in\unsure{\widehat{\mathbb{Z}}^{*}|\widehat{\mathbb{Z}}}$).
\deleted{Let} Thus by
definition, we have $\widetilde{u}\in\mathcal{M}_{0,3}^{\sim}$
(cf. \unsure{pages|pages} 390 ff).
Let $u'$ be the automorphism of $\mathcal{G}^{+\wedge}$ induced by $\widetilde{u}$
(via (3)) -- then $u'$ and $u$ induce the same
automorphism of $\mathcal{G}_{0,3}^{+\wedge}=\mathcal{G}^{+\wedge}/\{\pm1\}$,
thus we will have $u'=u\circ\Theta$, where
$\Theta$\inserted{$=u^{-1}u'$} is an automorphism which

% source: lower-middle portion, numbered properties and conclusion, lines 18--21.
$1^\circ)$ induces the identity on $\mathcal{G}^{+\wedge}/\{\pm1\}$, and

$2^\circ)$ takes $L_{0}^{\mathcal{G}}$ to a conjugate subgroup

-- we have seen that this implies $\Theta=1$, q.e.d.

% source: bottom portion, abandoned proposition and exact sequence (15).
\deleted{\uline{Proposition} Let $\widehat{\Pi}_{0}$ be the subgroup}
Let us now consider the exact \deleted{\unsure{lemma}}\inserted{sequence}

\diagram{\begin{tikzcd}
\text{(15)}\quad 1 \arrow[r] & \pm1 \arrow[r] & \widehat{\Pi} \arrow[r] \arrow[d,equal] & \Pi_{0,3}^{\wedge} \arrow[r] & 0\\
& & \mathcal{G}^{+\wedge}[2] & &
\end{tikzcd}}

\mspage{575}

% Printed archive footer outside the manuscript frame (not manuscript):
% 575
% https://grothendieck.umontpellier.fr
% Université de Montpellier, UMR CNRS 5149, Case courrier 051, Place Eugène Bataillon
% 34095 Montpellier cedex 5 - France

\scanpage{577}
\mspage{592}
% source: upper portion, lines 1--6, including (16)--(17).
and consider the splitting given by

\diagram{\begin{tikzcd}[column sep=large,row sep=large]
(16)\quad\left\{\vphantom{\begin{matrix}\Pi_{0,3}\\l_i\end{matrix}}\right.
& \Pi_{0,3} \arrow[r]
& \Pi=\mathcal{G}^{+}[2]=\operatorname{Ker}\bigl(\mathrm{Sl}(2,\mathbb Z)\longrightarrow\mathrm{Sl}(2,\mathbb Z/2\mathbb Z)\bigr) \\
& l_i \arrow[r,mapsto]
& \deleted{\unsure{\lambda_i|l_i}\overset{\unsure{\mathrm{dfn}}}{=}}\;\lambda_i\overset{\mathrm{dfn}}{=}-l_i
\end{tikzcd}}

There \deleted{\unsure{is}} is such a homomorphism, by virtue of
\[
(17)\qquad
\left\{\begin{aligned}
(-l_\infty)(-l_1)(-l_0)&=1\quad\text{in }\mathrm{Sl}(2,\mathbb Z)\quad\text{i.e.}\\
l_\infty l_1l_0&=-1
\end{aligned}\right.
\]
\margin{$)$}

and it is obvious that this is a splitting. We also likewise have a splitting for $\widehat\Pi$ over $\widehat\Pi_{0,3}$.

% source: middle portion, lines 7--12, definition (18).
We denote by $\Pi_0$ the image of $16)$, i.e.
\[
(18)\qquad
\left\{\begin{aligned}
\Pi_0&=\text{subgroup of }\Pi\inserted{=\mathcal{G}^{+}[2]}\text{ generated by}\\
&\lambda_0,\lambda_1,\lambda_\infty\qquad(\text{where }\lambda_i=-l_i)
\end{aligned}\right.
\]
and we thus have a subgroup $\widehat\Pi_0$ of $\widehat\Pi$.

% source: middle-lower portion, proposition a)--b), lines 13--18.
\uline{Proposition}\quad a) The subgroup $\Pi_0$ of $\Pi$ is \deleted{normal}\inserted{of index} $2$, and it is normal in $\mathcal{G}^{\unsure{\wedge|*}}$.

b) The subgroup $\widehat\Pi_0$ of $\widehat\Pi$ is \deleted{normal, \unsure{and}}\inserted{of index $2$}, and it is normal in $\mathcal M$ -- or, what amounts to the same, stable under the action of $\mathcal M/\pm1\simeq\mathcal M_{0,3}^{\sim}$.

% source: lower portion, proof and (19); right-hand additions at the same height.
\uline{Proof}\quad Since
\[
(19)\qquad
\left\{\begin{aligned}
\rho(\lambda_i)&=\lambda_{i+1}\\
\sigma(\lambda_0)&=\lambda_{1}^{\deleted{\unsure{-1}}},\quad
\sigma(\lambda_1)=\lambda_0^{\deleted{\unsure{-1}}}\\
\tau(\lambda_0)&=\lambda_0^{-1},\quad\tau(\lambda_1)=\lambda_1^{-1}
\end{aligned}\right.
\]
\margin{$\begin{aligned}
\sigma(\lambda_\infty)&=\sigma(\lambda_1\lambda_0)^{-1}=\lambda_1^{-1}\lambda_0^{-1}\\
\deleted{\unsure{\tau'}(\lambda_\infty)}&\deleted{=\lambda_\infty^{-1}}\\
\tau(\lambda_\infty)&=\tau(\lambda_1\lambda_0)^{-1}=\lambda_0\lambda_1
\end{aligned}$}

we see that $\Pi_0$ is \deleted{\unsure{stable}} normalized by $\rho,\sigma,\tau$, hence by $\mathcal G$. That it is of index $2$ is obvious. Thus $\widehat\Pi_0$ is of index $2$ in $\widehat\Pi$ and normalized by $\widehat{\mathcal G}$. To see that it is
\mspage{576}

% Printed archive footer outside the manuscript frame (not manuscript):
% 576
% https://grothendieck.umontpellier.fr
% Université de Montpellier, UMR CNRS 5149, Case courrier 051, Place Eugène Bataillon
% 34095 Montpellier cedex 5 - France

\scanpage{578}

% source-position: upper portion, continuation and display (20).
normalized by $\mathcal M$ in its \unsure{full} entirety, since
\[
\mathcal M=\mathcal M(0)\mathbin{.}\mathcal G^{+\wedge},
\]
we are reduced to verifying that it is normalized by
$\mathcal M(0)\simeq\mathcal M_{0,3}^{\sim}(0)$.
Now we \unsure{have}, for $u\in\mathcal M_{0,3}^{\sim}(0)$
\[
(20)\qquad
\left\{
\begin{aligned}
u(\rho)&=\operatorname{int}(\alpha)\rho\\
u(\sigma)&=\varepsilon_2\operatorname{int}(\beta)\sigma\\
u(\varepsilon_0)&=\varepsilon_0^\mu
\end{aligned}
\right.
\qquad
\left\{
\begin{aligned}
\varepsilon_2&=\varepsilon_2(\mu)\\
\alpha,\deleted{\unsure{\beta}}&\in\widehat{\Pi}\mathbin{.}\{1,\tau\}\\
\beta&\in\widehat{\Pi}
\end{aligned}
\right.
\]
% source-position: upper-middle portion, rewritten display (21).
\uline{whence}
\[
\deleted{u(\ell_0)=}\quad
\deleted{\unsure{(20)}}\quad
(21)\qquad
\left\{
\begin{aligned}
u(\ell_0)&=\ell_0^\mu\\
u(-\ell_0)&=(-\ell_0)^\mu
\quad\text{i.e.}\quad
u(\lambda_0)=\lambda_0^\mu
\quad(\text{since }(-1)^\mu=-1)
\end{aligned}
\right.
\]
\margin{$\unsure{)}$}

% source-position: middle portion, cancelled proof openings and lambda_1 calculation.
\deleted{It \unsure{remains} to \unsure{see}}
\deleted{\unsure{Indeed|Finally}}
thus $u(\lambda_0)\in\widehat{\Pi}_0$, it remains to prove that
$u(\lambda_1)\in\widehat{\Pi}_0$, \unsure{now}
\[
u(\lambda_1)
=u(\rho(\lambda_0))
=\underbrace{(u\rho)}_{\inserted{\deleted{\unsure{u}}(\operatorname{int}(\alpha)\rho)u}}(\lambda_0)
=\operatorname{int}(\alpha)(\rho)
\left(\underbrace{u(\lambda_0)}_{\inserted{\lambda_0^\mu}}\right)
=\deleted{\unsure{\operatorname{int}(\alpha)(\lambda_1^\mu)}}.
\]
\deleted{and we have \unsure{seen} that}
and since $\widehat{\Pi}_0$ is normalized by $\mathcal G^\wedge$
and in particular by $\deleted{\unsure{\rho}}\alpha\rho\alpha^{-1}$,
we indeed find $u(\lambda_1)\in\Pi_0$.

% source-position: lower portion, corollary (22).
\uline{Corollary} We have
\[
(22)\qquad
\left\{
\begin{aligned}
\Pi&\simeq\Pi_0\times\mu\\
\widehat{\Pi}&\simeq\widehat{\Pi}_0\times\mu
\end{aligned}
\right.
\]
where $\mu=\{\pm1\}=\operatorname{Centr}(\mathcal G^{+\wedge})$.

% source-position: bottom portion, cancelled and inserted proof of corollary.
\deleted{It suffices to \unsure{note} that the element $\unsure{-1}$
\unsure{of} $\mu$ \unsure{is not} in $\Pi_0$.}
\inserted{This is trivial, since we have a splitting of the central extension (15).}

\mspage{577}

% Printed archive footer outside the manuscript frame (not manuscript):
% 577
% https://grothendieck.umontpellier.fr
% Université de Montpellier, UMR CNRS 5149, Case courrier 051, Place Eugène Bataillon
% 34095 Montpellier cedex 5 - France

\scanpage{579}
\mspage{593}

% source-position: upper portion, introductory word and definitions (23)--(24).
\unsure{Let|Let}
\[
(23)\qquad
\left\{
\begin{aligned}
\mu_\tau&=\{1,\tau\}\subset\mathcal G
&D_\tau&=\mu_\tau\times\mu\subset\mathcal G\\
\mu_\sigma&=\{1,\unsure{\sigma|\tau'}\}\subset\mathcal G
&D_\sigma&=\mu_\sigma\times\mu\subset\mathcal G
\end{aligned}
\right.
\]
so that we have
\[
(24)\qquad
\begin{aligned}
D_\tau\cap\unsure{\mathcal G^+}
&=D_\tau\cap\Pi=\mu
&D_\tau\cap\Pi_0&=\{1\}\\
D_\sigma\cap\unsure{\mathcal G^+}
&=D_\sigma\cap\Pi=\mu
&D_\sigma\cap\Pi_0&=\{1\},
\end{aligned}
\]
\margin{\unsure{idem} on replacing $\mathcal G,\Pi$ by
$\mathcal G^\wedge,\widehat{\Pi}$}
% source-position: upper-middle portion, definitions (25)--(26).
We set
\[
(25)\qquad
\left\{
\begin{aligned}
\unsure{\Pi_\tau}&=\mu_\tau\mathbin{.}\Pi=D_\tau\mathbin{.}\Pi_0,
&\Pi_{0\tau}&=\mu_\tau\mathbin{.}\Pi_0\\
\Pi_\sigma&=\mu_\sigma\mathbin{.}\Pi=D_\sigma\mathbin{.}\Pi_0,
&\Pi_{0,\sigma}&=\mu_\sigma\mathbin{.}\Pi_0
\end{aligned}
\right.
\]
so that we have
\[
(26)\qquad
\left\{
\begin{aligned}
\Pi_\tau&=\Pi_{0\tau}\times\mu\\
\Pi_\sigma&=\Pi_{0,\sigma}\times\mu
\end{aligned}
\right.,
\]
and idem on putting in $\wedge$'s.

% source-position: middle-lower portion, normalization and (27).
\uline{Normalization of $\alpha$}, $\beta$\quad
The relations $20$, determine $\alpha,\beta$ only ``modulo sign''.
We remove this ambiguity by taking
\[
(27)\qquad
\alpha\in\Pi_{0,\tau},\qquad\beta\in\Pi_0.
\]

% source-position: lower portion, quotient presentation (28).
\uline{The group} $\mathcal G/\Pi_0$.
It is described by generators $\rho,\sigma,\tau^{\unsure{\prime}}$ and relations \unsure{by}
\[
(28)\qquad
\left\{
\begin{aligned}
&\rho^6=\sigma^4=\tau^2=1,\quad \rho^3=\sigma^2\\
&\tau(\rho)=\sigma(\rho^{-1}),\quad \tau(\sigma)=\sigma^{-1}\\
&\tau(\rho)=\rho
\end{aligned}
\right.
\]
% source-position: bottom portion, calculation (29) and revised concluding relations.
For we have
\[
(29)\qquad
\tau(\rho)=(-\ell_1^{-1})\rho
\]
thus the last relation $28$ reads
$\deleted{\tau(\rho)=\rho}\inserted{-\ell_1^{-1}=1}$
i.e. $\lambda_1=1$
(and implies
$\deleted{\unsure{-\ell_\infty=-\ell_0}}\inserted{\lambda_\infty=\lambda_0}=1$
on applying $\rho$).

\mspage{578}

% Printed archive footer outside the manuscript frame (not manuscript):
% 578
% https://grothendieck.umontpellier.fr
% Université de Montpellier, UMR CNRS 5149, Case courrier 051, Place Eugène Bataillon
% 34095 Montpellier cedex 5 - France